\subsection{Caracterização das Fontes de Dados}
A construção do Data Warehouse baseou-se na integração de múltiplas fontes de dados heterogéneas, com o objetivo de cruzar a informação operacional dos incêndios com variáveis meteorológicas, dados de observação terrestre e cobertura mediática.

\subsubsection{Registo de Ocorrências (ICNF)}
A fonte primária de dados consistiu numa versão tratada dos registos do repositório \texttt{icnf\_fire\_data}\footnote{\url{https://github.com/cityxdev/icnf_fire_data}}, que agrega o histórico de ocorrências do serviço \texttt{fogos.icnf.pt}.
Embora a fonte disponibilize registos desde 2001, este estudo incidiu sobre o período compreendido entre 2020 e 2025. Para garantir a relevância analítica, o \textit{dataset} foi submetido a um processo de filtragem rigoroso, resultando na exclusão de:
\begin{itemize}
	\item Registos com inconsistências geográficas (coordenadas em falta);
	\item Ocorrências com área ardida residual (inferior a 0,1 hectares);
	\item Registos incompletos, sem data de início ou fim de ocorrência.
\end{itemize}

\subsubsection{Dados Meteorológicos}
Para caracterizar o contexto atmosférico de cada ocorrência, recorreu-se à API do serviço \textit{Open-Meteo}\footnote{\url{https://open-meteo.com/}}. Através de um processo automatizado de extração, foram recolhidos dados históricos diários correspondentes às coordenadas geográficas e aos dias de atividade de cada incêndio, permitindo associar variáveis como temperatura, velocidade do vento e precipitação a cada evento específico.

\subsubsection{Contexto Mediático (Google News)}
Para avaliar o impacto mediático dos incêndios, foram recolhidos artigos noticiosos através de \textit{feeds} RSS do Google News. A extração foi filtrada por idioma (português) e por termos-chave associados a incêndios florestais no período em análise.
Estes dados não estruturados foram processados para extração de metadados (título, fonte, data), visando complementar a análise da atenção pública dada às ocorrências mais severas.

\subsubsection{Índices de Vegetação (MODIS)}
A monitorização do estado da vegetação foi realizada através do produto MOD13Q1 (NDVI de 16 dias com resolução de 250m), acedido via Google Earth Engine.
Para cada perímetro de incêndio, foram calculadas estatísticas zonais (média e contagem de pixeis) referentes aos períodos pré e pós-incêndio. Esta abordagem permitiu avaliar a biomassa disponível antes da ignição e quantificar o impacto na vegetação após a extinção.